\documentclass[UTF8,fontset=fandol,11pt,a4paper]{ctexart}
\usepackage[top=20mm,bottom=21mm,left=22mm,right=22mm,headheight=15pt]{geometry}
\usepackage{amsmath,booktabs,tabularx,array,graphicx,xcolor,caption,float}
\usepackage{fancyhdr,enumitem,setspace,tcolorbox,hyperref}
\definecolor{Navy}{HTML}{19354D}
\definecolor{Teal}{HTML}{267F79}
\definecolor{Muted}{HTML}{596979}
\definecolor{Pale}{HTML}{F1F6F6}
\definecolor{Line}{HTML}{DEE6EB}
\hypersetup{colorlinks=true,linkcolor=Navy,urlcolor=Teal,pdftitle={从交易闭环到真实样本：因子流动性回测阶段报告},pdfauthor={factor-liquidity 项目}}
\ctexset{section={format=\Large\sffamily\bfseries\color{Navy}},subsection={format=\normalsize\sffamily\bfseries\color{Navy}}}
\setstretch{1.17}
\setlength{\parskip}{0.45em}
\setlength{\parindent}{2em}
\setlist{leftmargin=1.6em,itemsep=0.3em,topsep=0.3em}
\captionsetup{font=small,labelfont=bf,labelsep=quad,skip=8pt}
\pagestyle{fancy}
\fancyhf{}
\fancyhead[L]{\small\sffamily\color{Muted} 因子流动性回测 · 阶段报告}
\fancyhead[R]{\small\color{Muted} 2026 年 10 月 14 日}
\fancyfoot[C]{\small\color{Muted}\thepage}
\renewcommand{\headrulewidth}{0.3pt}
\newcolumntype{Y}{>{\raggedright\arraybackslash}X}
\newcommand{\source}[1]{\par\noindent{\footnotesize\color{Muted}#1}\par}
\newcommand{\code}[1]{\texttt{\small\detokenize{#1}}}
\newenvironment{keybox}{\begin{tcolorbox}[colback=Pale,colframe=Line,boxrule=0.4pt,arc=1mm,left=3mm,right=3mm,top=2mm,bottom=2mm]}{\end{tcolorbox}}
\newcommand{\reportfigure}[3]{\begin{figure}[H]\centering\includegraphics[width=\linewidth]{figures/#1.pdf}\caption{#2}\label{#3}\end{figure}}
\begin{document}
\thispagestyle{empty}
\vspace*{9mm}
\noindent{\sffamily\color{Teal}\large FACTOR LIQUIDITY / RESEARCH NOTE}\par
\vspace{9mm}
\noindent{\sffamily\bfseries\color{Navy}\fontsize{30}{40}\selectfont 从交易闭环\par\noindent 到真实样本}\par
\vspace{4mm}
\noindent{\sffamily\color{Muted}\Large 因子流动性回测阶段报告}\par
\vspace{5mm}
\noindent{\color{Muted}合成数据验证、公司分钟行情回放、资金阶梯结果与后续研究}\par
\vspace{12mm}
\begin{keybox}
\textbf{本阶段已经完成的核心工作}\par
把目标权重转成可执行订单，接入参与率、部分成交、跨日余单、现金与持仓账本，跑通净值计算；随后接入真实的公司 1 分钟行情，以同一动量因子和组合规则完成资金阶梯回放。
\end{keybox}
\vspace{4mm}
\noindent\begin{tabularx}{\linewidth}{YYY}
\toprule
\textbf{163 项测试} & \textbf{50 个回测情景} & \textbf{6 组重绘图表}\\
当前代码全量通过 & 合成、真实各 25 个 & 汇总数据可离线重绘\\
\bottomrule
\end{tabularx}
\vspace{6mm}
\par\noindent
这份报告回答三个问题：已经做成了什么；图表能支持哪些判断；还需要补齐什么证据，才能讨论因子有效性和资金容量。

合成样本中，随着资金扩大，部分成交和余单逐渐增多。真实短样本中，策略净收益为负；参与率影响执行速度，但目前没有证据据此给出一个可部署的容量数字。

\textbf{两类结果承担不同任务。}合成数据用于验证工程与模型行为；真实数据用于观察一个固定历史区间内的条件性结果。两者不能放在同一收益排行榜里比较。
\vfill
\noindent{\small\color{Muted}报告日期：2026 年 10 月 14 日\quad{}｜\quad{}价格冲击模型：none\par
最低收益、最低完成率、最长延迟：尚未确定\par
本 bundle 包含正文、矢量图、绘图汇总数据和脚本；不包含公司原始分钟行情或连接凭据。}

\clearpage
\section{目前已经做完了什么}
项目已经从一个空仓库推进到可运行、可审计的回测原型。当前最有价值的成果，是能够把“因子想持有什么”与“账户实际买到了什么”分开记录，并逐笔解释差异。

\begin{table}[H]\centering\small
\begin{tabularx}{\linewidth}{p{26mm}YY}
\toprule
工作层次 & 已实现与已验证 & 证据的边界\\
\midrule
订单与账户 & 目标转订单、共享成交额度、冻结区间买入预算、部分成交、T+1 演示规则、费用与净值核对 & 证明程序遵守给定模型；不等于复刻真实撮合队列\\
合成实验 & 固定种子 42、手算案例、30 只模拟证券、80 个模拟日、25 个情景 & 检验资金和流动性约束是否作用于结果；不能证明因子有真实超额收益\\
真实数据接入 & 公司只读行情导出、字段检查、时间转换、竞价隔离、12 只证券、46 个交易日 & 行情是真实记录；历史状态、权益和实际可得时间尚未完整验收\\
真实阶梯回放 & 同一因子与组合规则、25 个情景、3,173 笔成交、全部账本审计通过 & 给出短样本、显式假设下的结果；不能直接宣布资金容量\\
\bottomrule
\end{tabularx}
\caption{已完成工作的层次及其证明范围}
\end{table}

\subsection{验证结果不是只看程序能否结束}
本次重新运行当前分支的全量测试，结果为 \textbf{163 passed}。手算入口也重新执行：20,000 股买入需求依次成交 10,000、5,000、5,000 股，最终现金 100,000 元、持仓 20,000 股、净资产 300,000 元，逐笔审计通过。

完整合成阶梯实际产生 76,004 笔成交，真实阶梯实际产生 3,173 笔成交；两个数据集各有 25 份审计记录，全部通过。审计包括成交前后现金和持仓、母子单关系、日子单费用、同股共享额度、每日持仓和净值。测试还覆盖账本被篡改时应当失败的反例，以及扰动未来数据不能改写既有决策的检查。

\subsection{建议怎样阅读后面的结果}
先看第 3 节的统一模型和第 4 节的合成压力实验，再看真实数据口径。真实结果要连续阅读净值、资金阶梯、完成率和时间四部分：单独看收益或平均完成时间，都容易漏掉订单所处的观察窗口。

\source{源代码验证分支为 codex/company-real-data。真实样本另在已提交的干净桌面副本上重跑，汇总数值与工作目录结果逐项相同。此报告整理并重绘既有回放结果，没有改变种子、证券池或执行参数来改善曲线。}

\clearpage
\section{两类数据分别回答什么问题}
合成数据给了我们一个可控制的实验环境：价格、成交量、交易状态和交易日由程序生成，可以构造共享额度、现金不足、部分成交和跨日余单等情形。因而，合成实验主要回答\textbf{实现是否正确、约束是否真正生效}。

真实数据把外部市场路径引入同一引擎，主要回答\textbf{在这些股票、这些日期和这些模型假设下，账户实际表现怎样}。它增加了现实相关性，却没有自动补齐历史交易制度、公司行动和因子统计检验。

\begin{table}[H]\centering\small
\begin{tabularx}{\linewidth}{p{31mm}YY}
\toprule
比较项目 & 合成样本 & 公司真实样本\\
\midrule
数据范围 & 30 只模拟证券，80 个模拟工作日 & 预先固定的 12 只深市证券，46 个真实交易日\\
行情粒度 & 5 分钟，48 根/模拟日 & 1 分钟；竞价有专门处理\\
源记录数量 & 115,200 行 & 133,032 行原始记录\\
区间 & 2025-01-02 至 2025-04-23 的模拟日历 & 实际观测 2026-07-03 至 2026-09-04\\
结果用途 & 程序正确性和模型压力响应 & 历史探索回放与数据接入验证\\
\bottomrule
\end{tabularx}
\caption{两类样本不是同一个市场的两个估计值}
\end{table}

\begin{keybox}
不能用“合成收益约 8.94\%，真实收益约 $-2.48$\%”来断言因子在现实中下降了多少。两者的价格形成方式、证券数量、行情粒度和样本长度都不同。合成收益的正负不承担因子有效性的证明责任。
\end{keybox}

\subsection{目前“因子验证”具体推进到了哪一步}
内置示例因子是 20 个交易日动量。已经证明，因子信号可以按照时间戳进入组合规则，组合规则可以进入执行和账户回放。真实样本还给出了这一规则在短区间内的净值与成交结果。

尚未完成的是：相对同股票池基准的超额收益、因子 IC/Rank IC、分组收益、不同市场阶段稳定性，以及权益修正和严格历史可得数据下的检验。只有 12 只候选证券而选择 Top 10，横截面筛选空间很小，不能把这一轮结果写成长期选股能力的结论。

\clearpage
\section{统一的组合与执行模型}
两类回放使用同一因子定义和同一组合规则。设 $P_{i,t}$ 为当日最后可见收盘代理价，则
\[
f_{i,t}=\frac{P_{i,t}}{P_{i,t-20}}-1.
\]
该定义需要 \textbf{21 个有效收盘价}。每 5 个交易日按动量排序，选择 Top 10 等权；收盘后形成信号，最早从下一可执行区间开始成交。相同规则不意味着每个执行情景的订单数量完全相同：费用、现金和先前成交会改变账户状态，之后目标股数随之变化。

\begin{table}[H]\centering\small
\begin{tabularx}{\linewidth}{p{34mm}Y}
\toprule
参数 & 本阶段取值\\
\midrule
初始资金 & 100 万、300 万、1,000 万、3,000 万、1 亿元\\
参与率 & 5\%、10\%、20\%\\
执行窗口 & 从首次可执行日计，最多 3 个交易日\\
调仓及现金缓冲 & 每 5 个交易日调仓；目标金额使用 0.5\% 现金缓冲\\
费用与数量规则 & 买入 0.03\%，卖出 0.08\%；每母单每日子单最低 5 元；100 股单位、T+1 可卖延迟均为显式演示规则\\
参考组 & 无参与率上限且不计费；无参与率上限且计费。两者仍保留现金、数量、可卖及数据约束\\
\bottomrule
\end{tabularx}
\caption{全部参数可配置；每笔资金对应 2 个参考组和 3 个受限组，共 25 个情景}
\end{table}

\subsection{成交额度和现金如何共享}
同一证券同一 bar 的所有订单共享 $pV_{i,b}$ 的成交额度，已成交数量在该额度中扣除；不会给每个订单重复分配一整份额度。订单只参与提交时间之后的区间，成交价采用该区间量额计算的 VWAP。VWAP 是回放的执行价格代理，不是下单时已知的价格，也不代表真实排队成交保证。

区间开始时冻结买入预算，所有买单共同消耗这份预算。区间内卖出所得从下一 bar 才可用于买入，费用一并占用现金。买入更新 $C\leftarrow C-qP-F$、$H\leftarrow H+q$；卖出更新 $C\leftarrow C+qP-F$、$H\leftarrow H-q$；净资产按 $NAV=C+\sum_i H_iP_i^{mark}$ 核对。

跨日延续的是母单剩余数量；新日期生成对应日子单。出现新目标时，先撤销旧余单，再按当前真实持仓重新计算差额，避免旧需求和新需求并存导致重复下单。窗口到期、目标替换与样本结束分别记录，不能统称为“流动性不足”。

\clearpage
\section{合成结果：约束已进入收益与成交路径}
\reportfigure{01_synthetic}{合成样本中 10\% 参与率的资金阶梯。收益与完成率使用原始回放汇总数值；参考组仍保留账户和交易模型约束。}{fig:synthetic}

\begin{table}[H]\centering\small
\begin{tabular}{rrr}
\toprule
初始资金 & 区间净收益 & 参考金额完成率\\
\midrule
100 万 & 8.9443\% & 100.0000\%\\
300 万 & 8.9136\% & 99.3510\%\\
1,000 万 & 8.8408\% & 92.8229\%\\
3,000 万 & 7.6125\% & 81.4444\%\\
1 亿元 & 6.0142\% & 69.0932\%\\
\bottomrule
\end{tabular}
\caption{合成样本的具体数值：80 个模拟日，非真实市场实证}
\end{table}

图 \ref{fig:synthetic} 显示，资金扩大后，订单相对于单根 bar 的成交额度变大，同一套规则需要更多时间完成交易。在 10\% 参与率下，完成率从 100\% 下降到约 69.09\%，净收益也从约 8.94\% 降到约 6.01\%。这说明订单没有被默认即时全成，持仓和净值确实随执行过程改变。

在 1 亿元时，提高参与率能够增加完成率：5\% 参与率约为 61.17\%，20\% 约为 77.06\%。但不能由此规定收益必须逐点单调。延迟成交会改变持仓时点，之后还会改变现金、订单和费用，实际收益路径可能交叉；本次没有调整种子去“修好”曲线。

\textbf{这一轮得到的是工程结论。}闭环对资金和成交约束有响应，账本能核对，因果时序按给定模型执行。合成市场中的正收益不构成真实 alpha，也不能用于给公司策略定容量。
\source{来源：固定种子 42 的 outputs/demo-final；完成率按母单决策参考价计量，完整定义见末节。}

\clearpage
\section{真实样本：接入了什么，仍缺什么}
本轮实际读取的是公司已有的 \code{stock_kdata.bar_1m}，没有用 mock 替代真实行情。证券池事先固定为 12 只深市股票，没有根据后续收益或成交额重新挑选。请求区间从 7 月 1 日开始，实际数据从 7 月 3 日开始；缺少的前两日没有补出行情。

\begin{table}[H]\centering\small
\begin{tabularx}{\linewidth}{p{38mm}Y}
\toprule
核对项 & 本轮观察与处理\\
\midrule
实际覆盖 & 46 个交易日，552 个股票日；每股票日 241 条源记录，共 133,032 行\\
时间戳 & 1 分钟表的 Z 标签是北京时间墙钟值的错误外观；去掉该标签后按北京时间解释，不再加 8 小时\\
数量与金额 & vol 为股，amount 为人民币元；均为分钟增量，不再做第二次差分\\
开盘竞价 & 09:30 共 552 条单独隔离，不并入首根连续交易 bar\\
回放与收盘窗 & 132,480 条进入回放；14:58、14:59、15:00 共 1,656 条只估值，不成交\\
量额价格差异 & 去除开盘竞价后，2,450 条 VWAP 位于快照 OHLC 之外；不修改源高低价来包围 VWAP\\
日历 & 使用深交所公布的 2026 年休市安排；尚待与公司交易日历及数据就绪记录交叉核对\\
\bottomrule
\end{tabularx}
\caption{公司 1 分钟样本的主要口径与检查结果}
\end{table}

VWAP 与 close 的绝对相对偏离最大约为 1.98\%。行情 OHLC 来源于 L2 快照聚合，而不是完整逐笔成交重建，因而不能把这个差异直接修成零，也不能据此证明真实队列的可成交性。观测区间的候选分钟网格完整，只说明这些股票日记录齐全，不代表全部历史证券状态已经认证。

\begin{keybox}
\textbf{真实价格 + 演示交易规则 + 市场事件时间回放。}\par
历史停复牌、ST、每日涨跌停价、数量和结算规则，以及公司行动账本尚未完整验收。15:00 使用源快照收盘代理；数据 D+1 批量可用且同键可修订，没有历史版本。本轮不是严格历史版本 PIT，也不是按数据库当时实际可得时间运行的生产策略回测。
\end{keybox}
\source{字段解释来自公司行情数据字典（2026-10-13 版本）并经实际接口和样本核对；公开 bundle 不收录内部连接信息。交易日历外部依据见参考资料。}

\clearpage
\section{真实净值：本区间没有显示正收益}
\reportfigure{02_real_nav}{真实样本的净值路径与预热期。资金 1 亿元；净值按初始资金归一化。阴影阶段用于积累 21 个收盘价，首个实际执行日为 8 月 3 日。}{fig:nav}

样本从 7 月 3 日开始，但 20 日动量需要 21 个有效收盘价。首次目标在 7 月 31 日周五 15:00 形成，首次成交从 8 月 3 日周一开始。46 个交易日中，前 21 日主要用于预热，实际持仓回放为后 25 个交易日。

以 100 万元为例，无参与率上限的毛收益为 $-2.4112$\%，计费后为 $-2.4820$\%；5\%、10\%、20\% 的结果与计费参考组相同。这说明在该小资金情景下，参与率并没有改变最终执行路径；亏损也不能全部解释成交易费用，因为不计费的参考组本身就是负收益。

图 \ref{fig:nav} 中，1 亿元计费参考组为 $-2.4848$\%，受限组约为 $-2.3760$\% 至 $-2.4235$\%。不同参与率曲线很接近，但路径和成交时点不完全相同。当前可以报告这个区间的负收益，不能据此断定动量因子长期无效，也不能声称低参与率改善了长期策略。

\subsection{平均现金占比为什么接近 46\%}
全区间平均包含了预热期的全现金账户，21/46 本身就约为 45.65\%。因此，约 46\% 的平均现金占比不能直接当作“买不进去”的证据。按回放的分钟净值记录取均值，只看 8 月 3 日起的持仓阶段，100 万、10\% 参与率平均现金约为 1.68\%；1 亿元、10\% 参与率约为 0.73\%。全区间和持仓阶段回答的是不同问题，比较时还必须保持采样频率一致。
\source{净值包含已计入模型的费用，不包含尚未接入的权益现金流。图中现金统计使用每分钟净值记录的等权均值；它不同于每日收盘均值，也不是盘中现金占用峰值。}

\clearpage
\section{真实资金阶梯：收益差异很小，不能据此定容量}
\reportfigure{03_real_returns}{真实样本的绝对净收益及相对计费参考组的差额。差额用基点表示，1 bp = 0.01 个百分点；正差额表示该路径少亏一些。}{fig:returns}

\begin{table}[H]\centering\small
\begin{tabular}{rrrr}
\toprule
资金 & 10\% 参与率净收益 & 最大回撤 & 模型费用/初始资金\\
\midrule
100 万 & $-2.4820$\% & 5.0488\% & 7.57 bp\\
300 万 & $-2.4787$\% & 5.0475\% & 6.96 bp\\
1,000 万 & $-2.4544$\% & 5.0123\% & 6.76 bp\\
3,000 万 & $-2.4563$\% & 4.9893\% & 6.75 bp\\
1 亿元 & $-2.3965$\% & 4.9818\% & 6.75 bp\\
\bottomrule
\end{tabular}
\caption{真实回放的资金阶梯：保持相同因子、组合规则及执行窗口}
\end{table}

资金增加后，净收益没有逐点下降。1 亿元时，5\%、10\%、20\% 参与率相对计费参考组分别高约 10.87、8.83、6.12 bp。这是本区间成交延迟与价格路径共同形成的结果，可能来自减少或推迟了某些亏损暴露；尚未做逐笔反事实分解，不能把这种解释写成已确认的机制。

账面费用可以直接核对，但“毛参考收益减净参考收益”不严格等于费用除以初始资金。计费会影响可买股数，之后现金、持仓与调仓订单也会改变。参考组用于比较完整路径，而不是把所有差额机械归入一个成本项。

\textbf{容量判断仍然没有触发条件。}最低收益、最低完成率、最长允许延迟均未确定，模型也未加入自身市场冲击。当前结果保留为条件性曲线，不输出“可承载 1 亿元”或“容量就是某档资金”的结论。

\clearpage
\section{完成率：必须把期末截断单独解释}
\reportfigure{04_real_completion}{真实样本完成率的两种观察口径。全部需求包括 9 月 4 日收盘新生成的订单；拥有完整可观察执行窗口的母单均完成。}{fig:completion}

10\% 参与率下，各资金的参考金额完成率约为 88.54\% 至 88.70\%。如果只看到这一列，很容易把其余约 11.4\% 理解成流动性不足。逐单检查后可以确认：\textbf{本次真实回放的所有期末未完成单，均在 9 月 4 日 15:00 提交，原因均为 sample\_end，尚未获得下一交易区间。}

所有情景中，已经拥有首次执行日且执行期限落在观测区间内的母单，数量完成率均为 100\%。这不是把不利订单删除后宣称成交良好，而是把两种事实并列展示：一类订单有完整观察机会；另一类在数据结束时刚产生，结果尚不可观察。

\begin{keybox}
100 万、10\% 参与率：53 个母单，44 个完成，9 个期末截断。\par
1 亿元、10\% 参与率：64 个母单，53 个完成，11 个期末截断。\par
“金额完成率约 88.6\%”与“完整观察母单全部完成”可以同时成立，因为分母不同。
\end{keybox}

\subsection{过程中确有约束，但不能与终态原因混为一谈}
执行事件中出现过参与率和现金限制，一些订单发生了部分成交；这些约束是真实参与计算的。只是本次最终残量都来自期末截断。最终未完成原因、盘中受限次数、完成时间和订单尾部，应该分别统计，不能把某个限制事件的出现等同于最终失败。

\subsection{下一轮应该怎样修正实验边界}
将信号截止日与执行观察截止日拆开：停止产生新目标后，继续提供至少一个完整执行窗口的行情和估值。报告同时保留全需求口径、完整观察口径和截断需求，不把最后一批订单当作容量瓶颈排序。本轮不借此改写原来的 88.6\% 数值，只增加解释和诊断。

\clearpage
\section{完成时间：平均自然小时掩盖了执行差异}
\reportfigure{05_real_clock}{真实订单的执行时间诊断。有效交易分钟从激活到末笔成交按可执行分钟网格累计；分位数和最大值按完整完成母单统计。}{fig:clock}

原汇总中的完成时间，从提交时刻计至最后一笔成交，包含隔夜、午休和周末。真实样本的调仓信号都落在周五收盘，下一次执行从周一 09:30 开始；这一段固定等待约为 \textbf{66.5 个自然小时}。因此，平均完成时间约 66 至 68 小时，并不表示连续交易了这么久。

为更接近执行问题，补充统计“模型允许执行的交易分钟”：从订单激活到最后成交，剔除午休、隔夜、周末和本模型禁止成交的收盘竞价分钟。本轮分钟网格完整，这个量反映订单跨过了多少可执行时间；它仍然不是券商实测排队耗时。

\begin{table}[H]\centering\small
\begin{tabular}{rrrrr}
\toprule
1 亿元参与率 & 平均有效分钟 & 中位数 & 90\% 分位 & 最大值\\
\midrule
5\% & 19.98 & 2 & 47.6 & 492\\
10\% & 9.74 & 2 & 25.2 & 226\\
20\% & 4.06 & 2 & 10.8 & 52\\
\bottomrule
\end{tabular}
\caption{相近的自然小时均值，背后存在明显不同的执行尾部}
\end{table}

100 万、10\% 参与率下，44 个完成母单平均只需约 1.20 个有效分钟，最大 2 分钟。到了 1 亿元，平均和尾部都拉长，尤其 5\% 参与率下最长母单用了 492 个有效分钟，跨越多个交易日。\textbf{时间尾部比全体平均值更早暴露流动性压力。}

参与率提高后，执行速度明显变快，但收益没有同样的单调规律。速度改善与投资结果改善是两件需要分别评估的事情，不能只用本区间的收益排序选择“最优参与率”。
\source{有效分钟按母单对应证券的 bar\_end > activated\_at 且 bar\_end $\leq$ last\_fill\_at 累计；均值为完成母单的等权平均，未按金额加权。}

\clearpage
\section{一笔跨日订单：余单怎样延续}
\reportfigure{06_order_timeline}{1 亿元情景的同一笔母单在三种参与率下的累计完成路径。横轴剔除非执行时段，并标出交易日边界；这是跨日余单的具体账本例子。}{fig:order}

以首批母单 O000004 为例：7 月 31 日 15:00 生成 000069.SZ 买入需求 5,653,400 股，决策参考价 1.76 元，参考金额约 995 万元；8 月 3 日 09:30 激活。三个情景的这笔初始需求相同，因此可以直接比较它的完成路径。

\begin{table}[H]\centering\small
\begin{tabular}{rrrr}
\toprule
参与率 & 最后一笔成交 & 有效交易分钟 & 提交至完成自然小时\\
\midrule
5\% & 8 月 5 日 09:48 & 492 & 114.80\\
10\% & 8 月 3 日 14:46 & 226 & 71.77\\
20\% & 8 月 3 日 10:22 & 52 & 67.37\\
\bottomrule
\end{tabular}
\caption{同一需求如何随参与率改变完成速度}
\end{table}

5\% 参与率下，8 月 3 日成交 3,003,500 股，8 月 4 日再成交 2,413,200 股，8 月 5 日成交剩余 236,700 股。三日合计正好为 5,653,400 股。母单延续剩余数量，每天的成交分别进入日子单和费用账本，最后在第 3 个执行日完成。

这一例子解释了为何“完整观察订单全部完成”仍可以与明显执行压力并存。订单没有失败，但低参与率已经把完成时间推到窗口第三日，资金在途中逐步变成持仓。若只看最终完成率，就看不到这段持仓偏离和时间风险。

图 \ref{fig:order} 是为了说明余单状态和时间尾部，\textbf{不代表这只股票可以据此确定容量}。执行采用历史量额 VWAP 代理，没有订单簿队列、价差、跨档和自身冲击；真实大额下单可能改变市场路径，本轮模型没有回答该问题。

\clearpage
\section{现在可以得出的结论}
\subsection{结论一：执行和账户闭环已经可以用来研究问题}
目标、订单、成交、费用、持仓和净值已能相互对应。固定种子、手算案例、真实接口回放和逐笔审计提供了不同层面的证据。下一轮可以替换因子或输入目标权重，沿用同一套执行与账本检查，而不必重新搭建基础闭环。

\subsection{结论二：合成实验已经显示资金约束会改变策略路径}
资金变大后，有限参与率使成交变慢、余单增加，实际持仓与目标逐步分离。合成样本中的收益和完成率变化支持这一模型响应。它给后续测试提供了一个受控基准，但不提供真实策略收益或容量的数字依据。

\subsection{结论三：当前真实短样本的投资结果为负}
本次示例动量组合在真实观测区间里没有显示正收益；不计费参考组也亏损。因此，不能把所有亏损归因于交易摩擦。另一方面，没有同池基准、权益修正和足够长的样本，也不能从这一点推出因子长期无效或没有超额收益。

\subsection{结论四：真实回放的完成率与时间需要分口径解释}
约 88.6\% 的金额完成率受末日新单截断影响；完整观察订单全部完成。约 66 小时的均值主要包含周末等待；大资金低参与率的交易时间尾部已经明显变长。这些诊断比一个平均数更能定位下一轮应观察什么。

\begin{table}[H]\centering\small
\begin{tabularx}{\linewidth}{YY}
\toprule
当前证据支持 & 当前证据不足以支持\\
\midrule
程序在给定模型下闭环且账本可核对 & 模型等价于真实市场成交能力\\
本区间真实净收益为负 & 动量因子长期无效或具有稳定 alpha\\
低参与率拉长部分大订单的完成时间 & 5\% 是最优参与率，或提高参与率必然改善收益\\
完整观察母单在当前窗口内完成 & 策略容量至少 1 亿元，或已经达到真实容量边界\\
\bottomrule
\end{tabularx}
\caption{把观察结果和外推结论分开}
\end{table}

\textbf{阶段定位：}本项目已完成工程闭环和真实小样本探索，接下来应完成数据与交易规则认证、修正观察边界，再进入因子研究。当前不设置业务阈值，不以“最高跑过的资金”代替容量。

\clearpage
\section{下一步：按证据缺口推进}
\subsection{第一优先级：把真实数据的关键口径验收完整}
核对公司交易日历、逐股逐日数据就绪和缺口；将源 15:00 快照与官方日收盘对照；接入历史停复牌、ST、每日涨跌停价、按日期生效的交易单位和可卖规则。公司行动不能只查“没有事件”：还应核对事件族覆盖状态、未解决记录，以及经济生效、到账和可交易日期。

这一阶段输出数据验收表和异常清单。对尚未认证的字段保留明确状态，不把缺记录默认解释为可交易。若要研究当时可部署的策略，还必须使用实际可得时间或历史版本；当前 D+1、可覆盖历史表的市场事件时间回放只能继续作为探索层。

\subsection{第二优先级：让订单观察窗口完整}
拆分信号生成区间和执行尾部，在最后一次信号后追加至少 3 个交易日的行情，停止生成新目标，继续执行和估值。固定“按提交时刻”与“按激活时刻”的两个时钟，同时输出完整观察、截断、到期、被替换和其他未成交原因；报告金额权重、订单等权及尾部分位数。

这一阶段应得到可比较的完成率和时间分布。是否能在追加样本后完成，是新实验的问题；不能提前把本轮截断订单当作已经成交。

\subsection{第三优先级：把示例因子变成可评估的研究对象}
先确定实际要研究的因子，保留信号时间戳和数据来源。增加同股票池等权基准、因子分组收益、IC/Rank IC、换手和风险暴露；把权益与现金流纳入账户或明确研究收益口径。扩大固定股票池和样本区间，覆盖不同市场阶段。高、中、低活跃度分层只能使用预热期已经可知的成交信息，不能按后续回测表现反选证券。

这一阶段输出基准对照、分组图和分时期结果，回答收益来自何处、是否稳定。新增图表必须依托实际计算，当前报告不预填这些结论。

\subsection{第四优先级：在业务目标明确后讨论容量}
由使用方确定最低收益、最低完成率与最长允许延迟，再用同一因子和规则扩展资金与参与率实验。先明确现有无冲击模型能支持的条件性容量范围；若要外推真实规模，还需要额外的执行价格和市场冲击证据。L2 跨档、自身冲击、复杂最优执行及 Web 界面仍不在本轮范围内。

\clearpage
\section{如何使用本报告 bundle 与已有回测代码}
\subsection{只看报告，或修改文字和图表}
公开仓库包含 \code{tex_bundle/report.pdf}、本正文 \code{main.tex}、6 组 PDF/PNG 图、绘图汇总 CSV 和 Python 脚本。已生成的 PDF 可直接阅读。修改正文后，在 \code{tex_bundle} 目录执行：
\begin{quote}\small\ttfamily
latexmk -lualatex -interaction=nonstopmode -halt-on-error \textbackslash{}\par
\hspace*{1em}-outdir=build main.tex
\end{quote}
\noindent 生成文件为 \code{build/main.pdf}；也可以用 \code{make pdf} 自动更新 \code{report.pdf}。所有图已经附在 bundle 中，编译正文无需数据库。使用 TeX Live / MacTeX 的 ctex 与 Fandol 字体即可。

需要改变颜色、线型或图表布局时，修改绘图脚本：\par
\noindent\code{scripts/make_figures.py}\par
\noindent 重绘命令为 \code{python scripts/make_figures.py}。该程序读取附带的汇总数据，可离线重绘，不会重新查询数据库或重新计算策略。

\subsection{完整回测应在已有引擎仓库运行}
公开报告仓库与已有回测引擎是两个不同交付物。完整的 orders、fills、events、positions、nav 与审计文件仍保留在本地引擎输出目录；公司原始分钟数据不随报告发布。进入已经安装依赖的引擎目录后，可执行：
\begin{quote}\small\ttfamily
PYTHONPATH=src .venv/bin/python -m pytest -q\par
PYTHONPATH=src .venv/bin/python -m factor\_liquidity.handcalc\par
PYTHONPATH=src .venv/bin/python -m factor\_liquidity.demo \textbackslash{}\par
\hspace*{1em}--config configs/demo.yaml --output outputs/my-demo\par
PYTHONPATH=src .venv/bin/python -m factor\_liquidity.company \textbackslash{}\par
\hspace*{1em}--config configs/company.example.yaml \textbackslash{}\par
\hspace*{1em}--output outputs/my-company-replay
\end{quote}
公司入口要求本地已有授权导出的文件及配置。每次使用新的输出目录。修改资金阶梯、参与率、窗口和费用使用 YAML 配置；若研究自己的因子，提供带 \code{decision_time, instrument, score} 的因子表，或带 \code{decision_time, instrument, weight} 的目标表，复用原文件入口。输入因子的历史可得性仍须由研究者提供证据。

\subsection{查看结果的顺序}
先读 \code{summary.csv} 看情景级表现，再查 \code{orders.csv} 的提交、激活、期限、终态和原因。成交、现金和费用看 \code{fills.csv}；每日子单看 \code{children.csv}；每日净值看 \code{daily_nav.csv}；审计看 \code{audit.json}。遇到结果差异时，先定位到具体母单和成交，再解释最终收益。

\clearpage
\section{指标定义、复现范围与参考资料}
\begin{table}[H]\centering\small
\begin{tabularx}{\linewidth}{p{33mm}Y}
\toprule
指标 & 本报告采用的定义\\
\midrule
区间净收益 & 末日净资产 / 初始资金 $-1$，包含模型费用；未处理公司行动的收益不称为完整总回报\\
最大回撤 & 净值相对此前峰值的最大下降比例\\
金额完成率 & $\sum_j q_j^{fill}P_j^{ref}/\sum_j q_j^{order}P_j^{ref}$；买卖均按正参考金额加总，不等于实际成交金额比例\\
完整观察母单 & 有首次执行日，且其执行期限不晚于观测截止日；本次这些真实母单均完全完成\\
自然完成小时 & 母单从 submitted\_at 至 last\_fill\_at 的自然小时，仅对完全完成母单统计\\
有效交易分钟 & 从激活到最后成交累计模型允许执行的分钟；剔除午休、隔夜、周末与禁成交竞价窗\\
平均现金比例 & 每分钟净值记录 cash/nav 的等权均值；“持仓阶段”在本样本中从 8 月 3 日起计算；日终均值另列，不能混用\\
参考净收益差额 & 受限组净收益 $-$ 相同资金计费参考组净收益；以 bp 显示，不能解释成单独的冲击成本\\
\bottomrule
\end{tabularx}
\caption{读图和跨情景比较应先统一指标定义}
\end{table}

\subsection{复现范围}
附带汇总数据能够重新生成本报告的六组图。\code{prepare_data.py} 能从已有完整回测输出重新提取这些汇总；它不会运行策略。完整回测需要本地引擎、原始行情和相应配置。报告采用固定参数，未为图形美观修改结果；重绘改变的是表达方式。

本次源回放的合成证据为 25 个情景、76,004 笔成交；真实证据为 25 个情景、3,173 笔成交。真实干净重跑对应引擎提交版本 \code{e0446eb}，汇总与工作目录结果逐项对照一致。报告整理过程中再运行全量测试和手算案例，结果见第 1 节。图表源数据、脚本说明和构建命令见 bundle README。

\subsection{资料来源}
\begin{enumerate}
\item 本地回测输出：固定种子合成阶梯、公司真实分钟回放，以及其逐笔账本、审计、配置与导出记录。本文数值以这些一手产物为依据。
\item 公司行情数据字典，2026-10-13 版本。内部资料，用于确认字段、单位、时间戳及快照聚合口径；公开 bundle 不包含内部连接信息。
\item 深交所：《关于 2026 年部分节假日休市安排的通知》，2025-12-22。\href{https://www.szse.cn/disclosure/notice/general/t20251222_618087.html}{深交所原文}。用于本次交易日历整理，不代替公司基础数据验收。
\item TDengine 官方 REST API 文档。\href{https://docs.tdengine.com/cloud/programming/client-libraries/rest-api/}{官方接口说明}。用于接口实现与返回格式核对。
\end{enumerate}

\end{document}
